\documentclass[11pt,a4paper]{article}
\usepackage[T1]{fontenc}
\usepackage[utf8]{inputenc}
\usepackage{lmodern,amsmath,amssymb}
\usepackage[margin=25mm]{geometry}
\usepackage{longtable,booktabs,array,calc}
\usepackage[hidelinks]{hyperref}
\hypersetup{pdftitle={One small-field blocking step, part 1b: the decay rate and the shape of the threshold, and},pdfauthor={navstokgap research working notes}}
\newcommand{\tightlist}{\setlength{\itemsep}{0pt}\setlength{\parskip}{0pt}}
\setlength{\emergencystretch}{3em}
\setcounter{secnumdepth}{0}
\begin{document}
\hypertarget{one-small-field-blocking-step-part-1b-the-decay-rate-and-the-shape-of-the-threshold-and-why-the-weak-sides-deficit-is-ten-orders-of-magnitude}{%
\section{One small-field blocking step, part 1b: the decay rate and the
shape of the threshold, and why the weak side's deficit is ten orders of
magnitude}\label{one-small-field-blocking-step-part-1b-the-decay-rate-and-the-shape-of-the-threshold-and-why-the-weak-sides-deficit-is-ten-orders-of-magnitude}}

Part 1 (\href{small-field-step-gaussian.md}{Gaussian level}) left one
constant unresolved, the exponential decay rate \(\kappa\) of the
fluctuation propagator, and part 2 was to produce the Kotecký--Preiss
threshold \(g_{\rm pert}^2\) of the remainder gas. This note settles
what can be settled about both without new mathematics, and the result
reshapes the weak-side programme. \textbf{The literature carries no
explicit rate.} Dimock's exposition of Balaban's construction
(Rev.~Math. Phys. 25 (2013) 1330010, §2.4; read from the arXiv text)
proves the decay by a parametrix built from Neumann inverses on cubes of
side \(M\) and a random-walk expansion that converges ``for \(M\)
sufficiently large'', with constants ``depending on \(L\) but no other
parameters''; the local estimate (80) there has a rate \(\gamma_0\)
stated only as \(O(L^{-2})\). \textbf{Explicit routes give
\(\kappa\sim10^{-3}\).} The generic Combes--Thomas argument gives
\(7\times10^{-4}\) for the hard constraint (part 1); the soft-constraint
operator \(-\Delta+aQ^TQ\) is local, has gap \(\tfrac29\) for
\(a\ge130\), and gives \(1.7\times10^{-3}\); the analytic-structure
route through the folded symbol is set up in Section 2 and reduces to a
zero-free strip for one explicit function, which is the defined
computation for an order-one rate. \textbf{The threshold scales as
\(\kappa^4\).} With the tree-graph organization of the remainder gas,
\[g_{\rm pert}\ \simeq\ \frac{1}{e\,C_3\,c_{\rm adj}\,\|C_{\rm fl}\|_{1,\kappa}/g^2}
\ \simeq\ \frac{\kappa^4}{2\times10^{4}\,C_3},\] with \(C_3\) the
cubic-vertex constant, of order \(10\) to \(30\), so that even at the
expected order-one rate \(\kappa\simeq\tfrac12\) the threshold is
\(g_{\rm pert}^2\sim10^{-12}\), and at the proven rate it is
\(10^{-35}\). Against the physical onset of the perturbative running at
\(g^2\simeq1\), \textbf{the weak side's explicit deficit is twelve
orders of magnitude in \(g^2\) with an ideal propagator and more than
thirty with the proven one, where the strong side's is a factor \(400\)
in \(\beta_W\).} The sub-target H1a of
\href{reasons-to-stop-as-research.md}{the research-directions note} is
therefore closed as a constant-chasing problem: no sharpening of the
present method reaches \(g^2\sim1\), and H1 is a methods problem. The
estimate of the threshold is a scaling estimate with explicit but crude
constants, labelled as such; the decay-rate statements are proved.
Constants explicit; nothing promoted.

\hypertarget{what-the-literature-does-and-the-soft-constraint}{%
\subsection{1. What the literature does, and the soft
constraint}\label{what-the-literature-does-and-the-soft-constraint}}

Dimock's Green's function is
\(G_k=(-\Delta+\bar\mu_k+a_kQ_k^TQ_k)^{-1}\): the block averaging enters
as a \textbf{soft} Gaussian constraint \(aQ^TQ\) rather than the hard
conditioning of part 1, and the operator is then local. The decay is
proved by covering the lattice with cubes of side \(M\), inverting with
Neumann conditions on each, gluing with a partition of unity, and
expanding the error in a random walk whose steps carry a factor
\(O(M^{-1})\); the walk converges for \(M\) large and the rate is half
the local rate. No numerical value is attached to \(M\), \(C\) or
\(\gamma_0\) anywhere in the construction, and Balaban's original
(Commun. Math. Phys. 95 (1984) 17) is of the same kind.

The soft constraint can be made explicit here. With \(Q\) the averaging
of part 1 (normalized by \(\tfrac1{16}\)),
\(QQ^T=\big(48+8(S+S^{-1})\big)/256\) along \(\mu\) has spectrum in
\([\tfrac18,\tfrac14]\), so \(\|Qf\|^2\ge\tfrac18\|f_\perp\|^2\) for the
component \(f_\perp\) of \(f\) orthogonal to \(K=\ker Q\). Writing
\(f=f_K+f_\perp\) and using
\(\|\partial f\|^2\ge\tfrac12\|\partial f_K\|^2-\|\partial f_\perp\|^2\)
with \(\|\partial f_K\|^2\ge\tfrac49\|f_K\|^2\) (part 1, Proposition 1)
and \(\|\partial f_\perp\|^2\le16\|f_\perp\|^2\),
\[\big\langle f,(-\Delta+aQ^TQ)f\big\rangle\ \ge\ \frac29\|f_K\|^2+\Big(\frac a8-16\Big)\|f_\perp\|^2
\ \ge\ \frac29\,\|f\|^2\qquad(a\ge130).\] The operator has range \(2\)
along \(\mu\) and \(1\) transversally, with \(\|Q^TQ\|\le\tfrac14\).
Conjugating by \(e^{\kappa x\cdot e}\) changes \(-\Delta\) by at most
\(2\sinh\kappa\) and \(aQ^TQ\) by at most \(\tfrac a4(e^{2\kappa}-1)\),
so the conjugated operator stays invertible with norm \(\le9\) when
\(2\sinh\kappa+\tfrac{130}4(e^{2\kappa}-1)\le\tfrac19\), that is
\[\kappa\ \le\ 1.7\times10^{-3},\qquad
\big|(-\Delta+aQ^TQ)^{-1}(x,y)\big|\le9\,e^{-\kappa|x-y|_\infty}.\] This
is a proved rate for the soft-constraint scheme, better than the
hard-constraint figure by a factor of two and of the same order: the
loss is intrinsic to any argument whose only inputs are the gap and the
hopping norm, since their ratio is about \(10^{-2}\).

\hypertarget{the-analytic-structure-route-set-up}{%
\subsection{2. The analytic-structure route, set
up}\label{the-analytic-structure-route-set-up}}

For the hard constraint, the fluctuation covariance is diagonal in the
coarse momentum \(\bar k\in[-\tfrac\pi2,\tfrac\pi2]^4\) and couples the
sixteen folded fine momenta \(k=\bar k+\pi n\), \(n\in\{0,1\}^4\). The
tube weight has the product form
\[|\hat w(k)|^2=\prod_{\nu\ne\mu}\cos^2\tfrac{k_\nu}2\ \cdot\ \cos^4\tfrac{k_\mu}2,
\qquad\hat k^2=\sum_\nu4\sin^2\tfrac{k_\nu}2,\] and folding exchanges
\(\cos^2\leftrightarrow\sin^2\) in the flipped components. With
\(s(\bar k)=\sum_n|\hat w(\bar k+\pi n)|^2/\hat k^2(\bar k+\pi n)\), the
diagonal entry of the fluctuation covariance at the unfolded momentum is
\[C_{\rm fl}(\bar k,\bar k)\ \propto\ \frac{\sum_{n\ne0}|\hat w(\bar k+\pi n)|^2/\hat k^2(\bar k+\pi n)}
{|\hat w(\bar k)|^2+\hat k^2(\bar k)\sum_{n\ne0}|\hat w(\bar k+\pi n)|^2/\hat k^2(\bar k+\pi n)},\]
in which the massless pole \(1/\hat k^2(\bar k)\) has cancelled
identically. Three facts fix where singularities can be. (a) For
\(\bar k\to\bar k+i\kappa e\) with \(\cosh\kappa<2\), every flipped
denominator satisfies
\(\operatorname{Re}\hat k^2(\bar k+\pi n)\ge2-2(\cosh\kappa-1)>0\), so
only the unfolded term can have a pole, and it has cancelled. (b)
\(|\hat w(\bar k)|^2\ge\tfrac1{32}\) throughout the strip, since each
factor has modulus at least \(\tfrac12\) there. (c) The remaining
singularities are the zeros of the denominator
\(D(\bar k)=|\hat w(\bar k)|^2+\hat k^2(\bar k)R(\bar k)\), with \(R\)
the flipped sum, analytic and bounded in the strip. \textbf{The decay
rate of the hard-constraint fluctuation propagator is the width of the
largest strip \(|\operatorname{Im}\bar k|<\kappa_*\) in which \(D\) has
no zero}, and \(\kappa_*\) is expected to be of order one. Establishing
it is a bounded computation with four real variables and one explicit
function, and it is the defined task for an order-one rate; no part of
it is done here.

\hypertarget{the-shape-of-the-threshold}{%
\subsection{3. The shape of the
threshold}\label{the-shape-of-the-threshold}}

Part 2 organizes the remainder \(S_W-S_2\), the Haar correction and the
non-abelian averaging as a gas of polymers on plaquettes, integrated
against the fluctuation Gaussian. Its structure is fixed by three
ingredients, whatever the details.

\emph{Vertices.} The cubic term of
\(N-\operatorname{Re}\operatorname{tr}U_p\) is
\(\operatorname{tr}\big(d\theta\,[\theta,\theta]\big)\)-type; with
fluctuations of root-mean-square size \(\tfrac32g\) per link (part 1),
\(d\xi\) is of size up to \(6g\) and the cubic vertex per plaquette is
of size
\(\tfrac2{g^2}\cdot6g\cdot\tfrac94g^2\cdot c_{\rm col}\simeq27c_{\rm col}\,g\),
so \(C_3g\) with \(C_3\) of order \(10\) to \(30\); quartic and Haar
terms are \(O(g^2)\) and subleading.

\emph{Covariance sums.} Factorizing the Gaussian integral of a product
over a set of plaquettes into connected polymers by the Brydges--Kennedy
interpolation puts one fluctuation covariance on each edge of a spanning
tree, so every polymer of \(n\) plaquettes carries a factor
\(\big(\|C_{\rm fl}\|_{1,\kappa}/g^2\big)^{n-1}\) with
\[\frac{\|C_{\rm fl}\|_{1,\kappa}}{g^2}=\sup_x\sum_y\frac{|C_{\rm fl}(x,y)|}{g^2}
\le\frac92\Big(\coth\frac\kappa2\Big)^4\ \simeq\ \frac{72}{\kappa^4}\qquad(\kappa\ll1),\]
using \(|C_{\rm fl}(x,y)|\le\tfrac92g^2e^{-\kappa|x-y|_\infty}\).

\emph{Entropy.} Plaquette adjacency \(20\) in four dimensions, so
connected sets of \(n\) plaquettes through a given one number at most
\((20e)^{n-1}\) (\href{wilson-strong-coupling-explicit.md}{Wilson note}
§2).

The Kotecký--Preiss condition then holds when the per-plaquette activity
times \(e\), the adjacency entropy and the covariance sum is below one,
\[e\cdot C_3g\cdot20e\cdot\frac{72}{\kappa^4}\ \le\ 1
\qquad\Longleftrightarrow\qquad
g\ \le\ g_{\rm pert}\simeq\frac{\kappa^4}{1.1\times10^4\,C_3},\] which
is the form quoted in the abstract.

\begin{longtable}[]{@{}
  >{\raggedright\arraybackslash}p{(\columnwidth - 6\tabcolsep) * \real{0.2500}}
  >{\raggedright\arraybackslash}p{(\columnwidth - 6\tabcolsep) * \real{0.2500}}
  >{\raggedright\arraybackslash}p{(\columnwidth - 6\tabcolsep) * \real{0.2500}}
  >{\raggedright\arraybackslash}p{(\columnwidth - 6\tabcolsep) * \real{0.2500}}@{}}
\toprule\noalign{}
\begin{minipage}[b]{\linewidth}\raggedright
\(\kappa\) (per lattice unit)
\end{minipage} & \begin{minipage}[b]{\linewidth}\raggedright
origin
\end{minipage} & \begin{minipage}[b]{\linewidth}\raggedright
\(g_{\rm pert}\) (\(C_3=20\))
\end{minipage} & \begin{minipage}[b]{\linewidth}\raggedright
\(g_{\rm pert}^2\)
\end{minipage} \\
\midrule\noalign{}
\endhead
\bottomrule\noalign{}
\endlastfoot
\(7\times10^{-4}\) & proved, hard constraint & \(10^{-18}\) &
\(10^{-36}\) \\
\(1.7\times10^{-3}\) & proved, soft constraint & \(4\times10^{-17}\) &
\(10^{-33}\) \\
\(0.5\) & expected order-one rate & \(3\times10^{-7}\) & \(10^{-13}\) \\
\(1\) & optimistic & \(5\times10^{-6}\) & \(2\times10^{-11}\) \\
\end{longtable}

This is a scaling estimate: the constants \(C_3\), the adjacency, and
the tree-graph organization are the standard ones and the form is
robust, but no polymer expansion has been carried out here, and a full
treatment adds factors of order one to ten in either direction. It does
not change the conclusion, because the conclusion is about orders of
magnitude.

\hypertarget{the-conclusion-for-the-weak-side}{%
\subsection{4. The conclusion for the weak
side}\label{the-conclusion-for-the-weak-side}}

The physical onset of the perturbative running is at \(g^2\simeq1\),
\(\beta_W\simeq6\), where asymptotic scaling is observed. The explicit
small-field step, done with the sharpest propagator rate one can hope
for and the standard organization, covers \(g^2\lesssim10^{-12}\); with
the proven rate, \(10^{-35}\). Every step between \(g^2_{\rm pert}\) and
\(g^2\simeq1\) is, in the rigorous sense, a step without a small
parameter, and their number is irrelevant since none is controlled. The
strong side's corresponding deficit, from the proved \(\beta_W=0.0135\)
to the physical \(5.7\), is a factor \(400\), and its methods are
polymer expansions of the same family. The difference is that the
strong-side gas has activities \(\beta_W/6\) per plaquette and a local
structure, while the weak-side gas has activities \(C_3g\) dressed by
covariance sums that scale like \(\kappa^{-4}\) and vertex constants of
order ten.

So H1a, the sharpness sub-target, is closed: sharpening constants within
the present method cannot bring \(g_{\rm pert}^2\) within twelve orders
of magnitude of the physical onset, and the meeting of the two reaches
through constant-chasing is not a viable route. H1 is a methods problem,
which is the content of the constructive programme extended through the
confinement scale: control of the renormalization steps at couplings
where the fluctuation is not small compared with the nonlinearity, with
no expansion in any parameter. That is the open problem, stated as
precisely as this programme can state it.

\hypertarget{consequence-for-state}{%
\subsection{5. Consequence for STATE}\label{consequence-for-state}}

The decay rate is proved at \(10^{-3}\) per lattice unit by two routes
and its order-one value is reduced to a zero-free strip for the explicit
function \(D\) of Section 2. The threshold of the remainder gas scales
as \(\kappa^4\) and lies at \(g^2\sim10^{-12}\) even for an ideal rate,
twelve orders below the physical onset of the running. The weak-side
sharpness sub-target is closed as not viable; the open problem is a
methods problem at order-one coupling, and the map of this programme
ends there.
\end{document}
